
Deep neural networks trained with empirical risk minimization (ERM) exploit spurious correlations present in training data. When a feature correlated with the label in training does not hold in deployment---or for minority subgroups within training---the model fails systematically. A classifier distinguishing waterbirds from landbirds may learn to rely on background (water vs. land) rather than bird morphology, because background provides a simpler, more consistent signal during training. When waterbirds appear on land---a minority of the training distribution---the model fails.

Existing solutions to this problem require either group labels during training (Group DRO; Sagawa et al., 2019) or two-stage procedures that first train a biased model, then identify failing samples for upweighting (JTT; Liu et al., 2021). While effective, these approaches cannot detect minority samples during a single training run. SPARE \citep{Yang2023} identifies spurious correlations early via simplicity bias but uses fixed epoch thresholds rather than per-sample trajectory analysis.

This work investigates whether per-sample loss trajectories contain discriminative information about minority group membership. The hypothesis is that simplicity bias---the tendency of neural networks to learn simpler patterns first---causes majority-group samples to achieve low loss rapidly via spurious features, while minority-group samples require core feature learning and exhibit persistently higher loss.

The contributions of this work are:

\begin{enumerate}
\item Demonstration that early-epoch loss magnitude provides a statistically significant signal for minority-group detection (Mann-Whitney p < $10^{-14})$ without requiring group labels.
\end{enumerate}

\begin{enumerate}
\item Verification via linear probes that simplicity bias manifests as a ~10 percentage point accuracy gap between spurious and core feature encoding throughout training.
\end{enumerate}

\begin{enumerate}
\item Explicit characterization of precision limits under minority imbalance, showing that the 5\% base rate---not signal weakness---bounds individual detection accuracy.
\end{enumerate}

\begin{enumerate}
\item A refinement of simplicity bias theory: with pretrained features, the bias manifests as an accuracy magnitude gap rather than a temporal gap.
\end{enumerate}

